\section{Riesgos materializados y Lecciones aprendidas}

\subsection{Riesgos materializados}

A los largo del proyecto, se han materializado algunos riesgos. Para nuestra suerte, la mayoria fueron previstos por el 
equipo. Estos riesgos previstos se trataron de cuestiones técnicas y falta de uso de algunas herramientas por parte del
equipo. Al inicio del proyecto, el equipo realizo una brainstorming y definió el alcance del proyecto. Como se trataba 
de una idea propia, ahi surgieron los primeros bocetos y funcionalidades que el equipo iba a realizar. 
\vspace{0.25cm}

Dada la alta motivación inicial, se presentaron muchas ideas que lamentablemente no vieron la luz. Creemos que esto se 
debe a una subestimacion de muchas de las tareas e inexperiencia en el uso de las herramientas como \textit{Unity}, por 
parte del equipo. A su vez hubo riesgos no previstos, quera cuestiones mas organizativas.
\vspace{0.25cm}

El equipo al inicio del proyecto, toma la determinación de desarrollar diferentes partes de forma aislada, dado que con la
idea de que cada parte tenga un \textit{owner} y un \textit{co-owner}. Por suerte, el tutor nos advirtió a tiempo sobre 
esta práctica y pudimos rectificarla. Sin embargo, en el inconsciente colectivo del equipo esta idea de mantener \textit{owner} 
y \textit{co-owner} en cada uno de los repositorios seguía presente y se materializo en la contribución a cada uno de ellos.
\vspace{0.25cm}

Las desventajas de separar cada repositorio entre un \textit{owner} y un \textit{co-owner}, es que desalienta el 
\textit{colective ownership} del proyecto completo y crea una necesidad de comunicación mas extendida entre los miembros 
del equipo. Pero dada la frecuencia de reuniones entre los miembros del equipo, el conocimiento de prácticamente todas 
las partes del código es casi total, para cada uno de los miembros del equipo.
\vspace{0.25cm}

Siguiendo con las cuestiones organizativas, el equipo tomo la decisión de poner las partes comunes entre el
\textit{game} y el \textit{world editor}, en un tercer repositorio llamado \textit{shared package}. A instancias finales 
del proyecto, los miembros del equipo encontraron esta decisión como algo tedioso y poco práctico. Es cierto que existian 
partes comunes entre el \textit{game} y el \textit{editor}, de esta forma se mitigarian el codigo repetido y se fomentaria 
las buenas prácticas de programación. 
\vspace{0.25cm}

Eso tornó poco practico el proceso de revisión y pruebas de las partes del código, ya que de realizar cambios solo o el 
\textit{game} o el \textit{editor}. No solo habría que probarlo en el repositorio afectado, sino que también en el otro. 
Añadiendo mas trabajo al proceso de revisión y a quien realizo los cambios, para que tenga que arreglar los \textit{bugs} 
introducidos.
\vspace{0.25cm}

Creemos que el trabajo hubiera sido mas facil sin la existencia de este tercer repositorio y se podría reemplazar con una
libreria propia solo para comunicación del \textit{core-service} y testeado contra el mismo o directamente sin este tercer 
repositorio.A su vez, no dudamos de la practicidad y versatilidad de la herramienta de control de versiones \textit{git}. 
Creemos que no es muy compatible con \textit{Unity}, dado que el mismo introducia archivos de metadatos donde estaban introducidos 
los datos propios de los objetos de \textit{Unity}. El equipo cometió algunas imprudencias propias de la metodología elegida 
de revisión de contribuciones, que llevo a conflictos mal resueltos, que luego se tradujeron en \textit{bugs}. O a sin querer 
omitir la subidad de estos archivos importantes.
\vspace{0.25cm}

En la primera parte del proyecto, el equipo se encontraba realizando este proyecto, pero a su vez cursando tres materias mas,
cada uno. Con lo cual, no hubo un desarrollo muy activo por parte del equipo, esto no solo se tradujo en un recorte, como 
mencionamos antes. Sino en un impacto un poco mas grande, poca comunicación, lo que se tradujo en poca cohesión entre los 
diferentes componentes del proyecto.
\vspace{0.25cm}

Esto es grave, por al conectar las diferentes partes del mismo, se presentaban \textit{bugs} y una desconexion total, lo 
que ponía en riesgo la realización del proyecto en sí. Este riesgo hubiese sido mitigado de forma total, si los miembros 
hubiésemos puesto en práctica el concepto del \textit{walking skeleton}. El mismo aparece mencionado en libro: 
\textit{Growing Object-Oriented Software, guided by Tests} de \textit{Steve Freeman} y \textit{Nat Pryce}, pero en realidad 
fue presentada por primera vez por \textit{Alistair Cockburn}, en su libro: \textit{Crystal Clear: A Human-Powered 
Methodology for Small Teams (Cockburn, 2005, p. 49)}.
\vspace{0.25cm}

El concepto \textit{walking skeleton}, de acuerdo con el primer libro, consiste en: \textit{"la implementacion mas chica 
posible, que pueda ser automaticamente construida, deployeada y testeada de punta a punta" (Freeman & Pryce, 2010, p. 32)}. 
Si bien este concepto, esta pensado para realizar entrega rápida y continua de producto y esta enmarcado en contexto de 
realizar fácilmente \textit{Test Driven-Development} (TDD). Esta práctica de realizar una implementacion mínima de punta 
a punta, hubiese facilitado y mitagado muchos riesgos de funcionalidad. Es el ejemplo de como buenas prácticas metodológicas 
y organizativas, pueden llevar a buenas prácticas funcionales.
\vspace{0.25cm}

Por otro lado, el equipo presento como un riesgo el colapso de la economía o la aparición de fraudes dentro del juego. Si
bien esto es un riesgo, y desde el equipo se presentó la posibilidad de realizar un modelo económico. Por el riesgo mayor
que presentaba las decisiones arquitectónicas y la falta de cohesion entre los diferentes componentes del proyecto, se 
decidió darle mas prioridad a este riesgo, que al modelo económico. Finalmente, se recorto el alcance del mismo a un modelo 
económico-financiero para analizar la viabilidad del proyecto y la puesta en marcha del proyecto, a través de
inversiones semilla o inversores ángeles. Es decir, que la parte de modelado, análisis y detección de fraude queda para 
trabajo futuro.
\vspace{0.25cm}

\subsection{Lecciones aprendidas}

A continuación mencionamos cuales fueron las lecciones aprendidas por todos los miembros del equipo:
\vspace{0.25cm}

\textbf{\underline{Comunicación}}: Si bien no hubo mucha interacción entre los miembros del equipo, en la primera parte 
del proyecto y muchos presentaron agotamiento y/o cansancion de mantener reuniones 3 veces por semana. La comunicación
resultó ser el pilar fundamental para sobrellevar el proyecto en la segunda parte del mismo, estan no solo mantuvo la
moral alta del equipo en momentos de tensión. Permitió el conocimiento total de cada parte del proyecto por cada miembro,
compartir y delegar tareas entre todos los miembros. Es importante destacar que la comunicación no solo fue mantenida cara
a cara, sino que a través de chats de \textit{WhatsApp} de forma grupal y uno a uno entre todos.
\vspace{0.25cm}

\textbf{\underline{Planificación, organización y coordinación}}: El equipo pudo llevar el estado al proyecto actual, 
a pesar de no estar muy activo en la primera parte. Gracias una correcta planificación de cada una de las decisiones que 
se fueron tomando, las decisiones no fueron solo de implementacion, sino también de diseño. El hecho de que el equipo se haya
organizado y coordinado a través de herramientas colaborativas como \textit{Miro} o \textit{Excalidraw}, no solo para pensar 
en como resolver el problema, sino para preever casos borde y posibles deficiencias en cada decision, llevo a lograr de 
forma exitosa el MVP, a pesar de los contratiempos.
\vspace{0.25cm}

\textbf{\underline{Sobreestimacion y subestimacion de las tareas}}: El hecho de que la mayoria de las tareas haya sido 
subestimadas y algunas de ellas sobreestimadas. Llevo a varios recortes de funcionalidades y de alcances del mismo. 
Esto se debe a que, como dijimos antes, el equipo estaba altamente motivado por el proyecto. Sin embargo, no esperaba 
las consecuencias de una reducción de actividad y un alto aislamiento en la primera parte del proyecto. Esto nos lleva a 
pensar, que el proceso creativo es importante, pero tambien en que hacer todo no es posible y hubiese sido necesario revisar 
de forma mas frecuente cada una de las funcionalidades.
\vspace{0.25cm}

\textbf{\underline{Impacto de las decisiones de diseño e implementacion}}: Las decisiones son el medidor del éxito de un 
proyecto, aqui aparece un contraste interesante. Desde el equipo, planificamos casi perfectamente las decisiones de diseño 
mitigando casos borde y posibles \textit{bugs}. Pero a la hora de la verdad, es decir la implementacion, encontramos dificultades 
por no saber como desarrollar ciertas partes, como el editor de mundos o el despliegue en servidores. Esto nos llevó a 
refactotizar alguna partes del código, algunas incluso varias veces. Algunas fueron resueltas de manera colectiva, otras 
gracias a nuestro tutor. Creemos que lo importante Para mitigar esto en futuros proyectos, es a través de colaboración 
colectiva, un aprendizaje mayor de la herramientas e investigación de las mismas, antes de un dirigirse directamente a 
implementar.
\vspace{0.25cm}
